\section{ADCS Threat Model and Systematic Blind Spots}
\label{sec:threat_model}

In this section, we formalize the heterogeneous identity multigraph of Active Directory Certificate Services and expose two fundamental blind spots in current defensive audit practices: the static snapshot blind spot under rollback attacks, and the PKCS\#10 ASN.1 schema discontinuity between CA and KDC telemetry.

\subsection{Heterogeneous Access Control Multigraph Formalization}
We model an enterprise Active Directory forest and its integrated Public Key Infrastructure as a directed heterogeneous multigraph.

\begin{definition}[Identity Multigraph Schema]
Let $\mathcal{G} = (\mathcal{V}, \mathcal{E}, \phi, \psi)$ be a directed heterogeneous multigraph, where:
\begin{itemize}
    \item $\mathcal{V}$ is the finite set of identity and PKI vertices, and $\phi: \mathcal{V} \to \mathcal{T}_V$ is a node-type assignment function mapping vertices to the entity set:
    \begin{equation}
    \mathcal{T}_V = \{\text{User}, \text{Group}, \text{Computer}, \text{CertificateTemplate}, \text{EnterpriseCA}, \text{Domain}\}
    \end{equation}
    \item $\mathcal{E}$ is the finite set of directed edges, and $\psi: \mathcal{E} \to \mathcal{T}_E$ maps edges to relationship types:
    \begin{align}
    \mathcal{T}_E = \{&\text{MemberOf}, \text{GenericAll}, \text{GenericWrite}, \text{WriteOwner}, \text{WriteDacl}, \nonumber \\
    &\text{Enroll}, \text{ExtendedRight}, \text{EnablesEnrollment}, \text{IssuedBy}, \text{TrustedForNTAuth}\}
    \end{align}
    \item Each node $v \in \mathcal{V}$ is associated with a raw feature vector $\mathbf{x}_v \in \mathbb{R}^{d_{\phi(v)}}$. For certificate templates $t \in \mathcal{V}_{\text{Template}}$, $\mathbf{x}_t \in \{0, 1\}^{10}$ encodes security configuration flags, including whether the template permits enrollee-supplied SANs (\texttt{CT\_FLAG\_ENROLLEE\_SUPPLIES\_SUBJECT}), requires manager approval (\texttt{CT\_FLAG\_PEND\_ALL\_REQUESTS}), or specifies Client Authentication EKUs (\texttt{1.3.6.1.5.5.7.3.2}).
\end{itemize}
\end{definition}

An attack path $p$ from an initial compromised user foothold $u_{\text{attacker}} \in \mathcal{V}_{\text{User}}$ to a high-value domain target $v_{\text{target}} \in \mathcal{V}_{\text{Domain}}$ is a sequence of alternating vertices and edges:
\begin{equation}
p = (v_0, e_1, v_1, e_2, \dots, e_k, v_k), \quad v_0 = u_{\text{attacker}}, \quad v_k = v_{\text{target}}
\end{equation}
In an ADCS privilege escalation, the path traverses intermediate membership and ACL delegation edges to a certificate template $t$, leverages an enterprise CA $c$ to issue an administrative credential, and culminates in a domain-level credential takeover.

\subsection{The Static Snapshot Blind Spot: Formal Proof and Empirical Evidence}
Existing enterprise defense workflows rely on periodic directory extraction tools (e.g., BloodHound, SharpHound, Certipy) that execute every 6 to 24 hours. Under temporal mutation-and-rollback attacks (such as ESC4 and ESC9), an attacker exploits transient write permissions to mutate a template, obtain a rogue certificate, and immediately roll back the attributes to their original configuration.

\begin{definition}[Static Snapshot Invariance under Rollback]
Let $t_0 < t_1 < t_2 < t_3$ denote the sequence of discrete timestamps in an attack execution lifecycle:
\begin{itemize}
    \item $t_0$: Pristine baseline state of directory $\mathcal{G}_{t_0}$.
    \item $t_1$: The attacker executes an LDAP modification adding \texttt{Enroll} permissions or setting \texttt{CT\_FLAG\_ENROLLEE\_SUPPLIES\_SUBJECT} on template $t$, yielding state $\mathcal{G}_{t_1}$.
    \item $t_2$: The attacker submits a certificate request and receives an issued certificate $\text{Cert}_{\text{admin}}$ valid for 1 year.
    \item $t_3$: The attacker rolls back the template modifications, restoring the directory attributes to their baseline state, yielding $\mathcal{G}_{t_3}$.
\end{itemize}
\end{definition}

\paragraph{Empirical Verification on Live Windows Server 2019:}
To measure the impact of this vulnerability, we provisioned an enterprise ADCS cyber-range on Windows Server 2019 Domain Controllers and Enterprise CAs. At each phase of an automated ESC4 attack executed by an unprivileged domain account (\texttt{ESSOS\textbackslash khal.drogo}), we extracted the entire Active Directory database and computed cryptographic SHA-256 digests over the canonicalized JSON object graph. The resulting hashes were:
\begin{align}
\text{SHA-256}(\mathcal{G}_{t_0}) &= \texttt{219376ed1562fb8471f25af34b02d0f18cb044e8ebfe0be907c21e1ddf6452d8} \\
\text{SHA-256}(\mathcal{G}_{t_1}) &= \texttt{d4f82a991820bc695361de42ab5a782806ccbf269a8bfe7df2f518e95079a491} \\
\text{SHA-256}(\mathcal{G}_{t_3}) &= \texttt{219376ed1562fb8471f25af34b02d0f18cb044e8ebfe0be907c21e1ddf6452d8}
\end{align}

Because $\text{SHA-256}(\mathcal{G}_{t_0}) = \text{SHA-256}(\mathcal{G}_{t_3})$, the directory state at $t_3$ is bit-identical to the pre-attack state. Any security system that evaluates static snapshots at scheduled intervals $T > t_3 - t_0$ experiences a complete detection failure:
\begin{equation}
\text{Recall}_{\text{static}}(\text{Attack}_{t \ge t_3}) = 0.0\%
\end{equation}
The attacker retains the cryptographic certificate $\text{Cert}_{\text{admin}}$ indefinitely, enabling Kerberos PKINIT authentication at arbitrary future timestamps without leaving any residual footprint in the directory multigraph.

\subsection{Protocol Discontinuity: PKCS\#10 ASN.1 SAN Concealment and CA--KDC Logging Gap}
Even when defenders deploy Security Information and Event Management (SIEM) systems to collect Windows Security audit logs, standard detection rules fail due to a structural discontinuity between Certificate Authority and Key Distribution Center logging schemas.

\begin{proposition}[CA--KDC Schema Disjointness]
Let $\pi_{\text{CA}}$ denote the projection of security attributes recorded by Microsoft-Windows-CertificationAuthority (Event IDs 4886 and 4887), and let $\pi_{\text{KDC}}$ denote the projection recorded by Microsoft-Windows-Security-Auditing during Kerberos PKINIT pre-authentication (Event ID 4768). Then:
\begin{equation}
\pi_{\text{CA}} \cap \pi_{\text{KDC}} = \emptyset
\end{equation}
Specifically, $\pi_{\text{CA}} = \{\text{Subject}, \text{TemplateName}, \text{RequestID}, \text{Disposition}\}$, whereas $\pi_{\text{KDC}} = \{\text{TargetUserName}, \text{PreAuthType}=16, \text{CertSerialNumber}, \text{CertThumbprint}\}$.
\end{proposition}

\paragraph{The ASN.1 SAN Concealment Phenomenon:}
When an attacker requests a certificate with a rogue Subject Alternative Name (e.g., impersonating the Domain Administrator), the SAN attribute can be supplied through two distinct mechanisms:
\begin{enumerate}
    \item \textbf{Explicit RPC Attribute String:} Passed via the \texttt{Attributes} string parameter of \texttt{ICertRequest::Submit} (e.g., \texttt{"SAN:upn=Administrator@domain.local"}). In this configuration, the CA security log explicitly populates the \texttt{Attributes} field in Event ID 4887.
    \item \textbf{Raw PKCS\#10 ASN.1 Binary Extension:} Embedded directly within the cryptographic Certificate Signing Request (CSR) under the \texttt{szOID\_SUBJECT\_ALT\_NAME2} (\texttt{2.5.29.17}) extension attribute.
\end{enumerate}

Our forensic investigation of live Windows Server 2019 event streams reveals that when an attacker employs the PKCS\#10 ASN.1 CSR extension (as executed natively by modern exploit tooling), the CA event log omits the SAN string entirely from the XML event payload:
\begin{lstlisting}[language=XML,caption={Live Windows Server 2019 Event ID 4887 under PKCS\#10 ASN.1 SAN Concealment.},label=lst:ca_event]
<Event xmlns="http://schemas.microsoft.com/win/2004/08/events/event">
  <System>
    <Provider Name="Microsoft-Windows-CertificationAuthority"/>
    <EventID>4887</EventID>
    <TimeCreated SystemTime="2026-10-17T03:15:22.418Z"/>
  </System>
  <EventData>
    <Data Name="RequestId">19</Data>
    <Data Name="Requester">ESSOS\khal.drogo</Data>
    <Data Name="Template">ESC1-Vulnerable</Data>
    <Data Name="Subject">khal.drogo</Data>
    <Data Name="Attributes"></Data> <!-- CRITICAL: SAN is completely absent! -->
  </EventData>
</Event>
\end{lstlisting}

As demonstrated in Listing~\ref{lst:ca_event}, the \texttt{Attributes} tag is completely blank, and the \texttt{Subject} field displays only the low-privileged requester's identity (\texttt{khal.drogo}). Consequently, any SIEM correlation rule that inspects Event 4887 for unauthorized administrator UPNs fails with a $100\%$ false-negative rate. When the attacker subsequently authenticates as \texttt{Administrator} via Kerberos PKINIT (Event 4768), the KDC log records only the certificate serial number and SHA-1 thumbprint, leaving zero common identifiers to link the issuance to the compromised authentication.
